\chapter*{DANH MỤC KÝ HIỆU VÀ TỪ VIẾT TẮT}
\addcontentsline{toc}{chapter}{Danh mục ký hiệu và từ viết tắt}

\section*{1. Ký hiệu toán học và biến trạng thái}

\begin{table}[h!]
\centering
\begin{tabularx}{\textwidth}{l l X}
\toprule
\textbf{Ký hiệu} & \textbf{Đơn vị / Miền giá trị} & \textbf{Ý nghĩa sinh học / toán học} \\
\midrule
$t$ & Năm (Years) & Biến thời gian trong mô hình động thái quần thể \\
$x(t)$ & $10^9$ cá thể & Sinh khối / Mật độ quần thể cá mòi dầu (\textit{Atlantic menhaden} -- con mồi) \\
$y(t)$ & $10^9$ cá thể & Sinh khối / Mật độ quần thể cá vược sọc (\textit{Striped bass} -- động vật săn mồi) \\
$r$ & $\text{năm}^{-1}$ & Tốc độ tăng trưởng nội tại của quần thể con mồi \\
$a$ & $\text{năm}^{-1} \times (\text{tỷ con})^{-1}$ & Hệ số tự giới hạn (cạnh tranh nội bộ loài) của con mồi \\
$b$ & $\text{năm}^{-1} \times (\text{tỷ con})^{-1}$ & Tỷ lệ cá mòi dầu bị cá vược sọc săn bắt (tác động của kẻ săn mồi) \\
$e$ & $\text{năm}^{-1}$ & Tỷ lệ tử vong tự nhiên của cá vược sọc khi không có con mồi \\
$c$ & $\text{năm}^{-1} \times (\text{tỷ con})^{-1}$ & Hệ số tự giới hạn (cạnh tranh nội bộ loài) của kẻ săn mồi \\
$d$ & $\text{năm}^{-1} \times (\text{tỷ con})^{-1}$ & Hiệu suất chuyển đổi dinh dưỡng thức ăn thành sinh khối cá vược sọc \\
$E_1, E_2$ & Không thứ nguyên & Nỗ lực khai thác đánh bắt đối với cá mòi dầu và cá vược sọc \\
$q_1, q_2$ & $[0, 1]$ & Hệ số khả năng đánh bắt (\textit{catchability coefficients}) của từng loài \\
$p_1, p_2$ & Đơn vị tiền tệ / sinh khối & Giá bán trên một đơn vị sinh khối của từng loài \\
$c_1, c_2$ & Đơn vị tiền tệ / nỗ lực & Chi phí vận hành trên một đơn vị nỗ lực khai thác của từng loài \\
$\pi(E_1, E_2)$ & Đơn vị tiền tệ / năm & Hàm tô kinh tế (lợi nhuận ròng hàng năm từ khai thác) \\
$\tau$ & Năm & Thời gian hồi phục của hệ sinh thái về trạng thái cân bằng \\
$R$ & $\text{năm}^{-1}$ & Độ đàn hồi / độ hồi phục sinh học (\textit{Resilience}), $R = 1/\tau = -\mathrm{Re}(\lambda_m)$ \\
$\Omega$ & Miền không gian pha & Miền bất biến dương sinh học của hệ phương trình vi phân \\
$S$ & Không thứ nguyên & Sai số chuẩn phần dư trong quá trình hồi quy phi tuyến \\
\bottomrule
\end{tabularx}
\end{table}

\section*{2. Các từ viết tắt chuyên ngành}

\begin{table}[h!]
\centering
\begin{tabularx}{\textwidth}{l l X}
\toprule
\textbf{Từ viết tắt} & \textbf{Tên tiếng Anh} & \textbf{Ý nghĩa tiếng Việt} \\
\midrule
\textbf{MSY} & Maximum Sustainable Yield & Sản lượng bền vững tối đa \\
\textbf{MEY} & Maximum Economic Yield & Sản lượng kinh tế tối đa \\
\textbf{RMY} & Resilience Maximizing Yield & Sản lượng tối đa hóa độ đàn hồi sinh thái \\
\textbf{CPUE} & Catch Per Unit Effort & Sản lượng đánh bắt trên một đơn vị nỗ lực \\
\textbf{EBFM} & Ecosystem-Based Fisheries Management & Quản lý nghề cá dựa trên hệ sinh thái \\
\textbf{ERP} & Ecological Reference Point & Điểm tham chiếu sinh thái \\
\textbf{ASMFC} & Atlantic States Marine Fisheries Commission & Ủy ban Nghề cá Biển các Bang Đại Tây Dương \\
\textbf{SEDAR} & Southeast Data, Assessment, and Review & Cơ quan Đánh giá Dữ liệu Trữ lượng Đông Nam Hoa Kỳ \\
\textbf{NEFSC} & Northeast Fisheries Science Center & Trung tâm Khoa học Nghề cá Vùng Đông Bắc Hoa Kỳ \\
\textbf{NOAA} & National Oceanic and Atmospheric Admin. & Cơ quan Quản lý Khí quyển và Đại dương Quốc gia Mỹ \\
\textbf{CBF} & Chesapeake Bay Foundation & Quỹ Bảo tồn Vịnh Chesapeake \\
\textbf{GDP} & Gross Domestic Product & Tổng sản phẩm quốc nội \\
\bottomrule
\end{tabularx}
\end{table}
